\chapter{Formatos de log soportados}
\label{sec:log-formats}

\section{El trait \code{LogParser}}

El punto de extensión está en
\filepath{platform/parsers/src/lib.rs}:

\begin{lstlisting}[language=Rust]
pub trait LogParser: Send + Sync {
    fn name(&self) -> &str;
    fn can_parse(&self, sample: &[u8]) -> bool;
    fn parse(&self, bytes: &[u8]) -> Result<Vec<Event>>;
    fn field_roles(&self) -> &FieldConfig;
}
\end{lstlisting}

El dispatcher \code{format_registry} en
\filepath{platform/api/src/format_registry.rs} consulta
\code{can_parse} sobre una muestra del archivo y elige el primer
parser que responde positivamente. El fallback es el parser
\code{generic} con heurísticas.

Cada parser produce un \code{Vec<Event>}; cada \code{Event} es una
tupla \code{(timestamp, actor, action, target, attrs)} con roles
de campo ya asignados.

\section{Sysmon (Windows Event Log XML)}

\textbf{Fuente}: archivos EVTX exportados, o XML de
\code{Microsoft-Windows-Sysmon/Operational}. Parser:
\filepath{platform/parsers/src/sysmon.rs} + lógica antigua en
\filepath{core/graph-engine/src/sysmon.rs}.

\textbf{Event IDs cubiertos}: 1 (process create), 3 (network
connect), 7 (image load), 8 (CreateRemoteThread), 10
(ProcessAccess), 11 (FileCreate), 13 (RegistryValueSet), 22 (DNS),
23 (FileDelete).

\textbf{Mapeo a entidades}:
\begin{itemize}
    \item \code{Image}, \code{ParentImage} $\to$ \code{Process}.
    \item \code{User} $\to$ \code{User}.
    \item \code{Computer} $\to$ \code{Host}.
    \item \code{DestinationIp} $\to$ \code{IP}.
    \item \code{QueryName} $\to$ \code{Domain}.
    \item \code{TargetFilename} $\to$ \code{File}.
\end{itemize}

\textbf{Relaciones típicas}: \code{Process -[Spawn]-> Process}
(ID 1), \code{Host -[Connect]-> IP} (ID 3),
\code{Process -[Write]-> File} (ID 11).

\section{Sentinel / Log Analytics (JSON)}

\textbf{Fuente}: resultados de query KQL exportados como JSON o
streamed via conector. Parser:
\filepath{platform/parsers/src/sentinel.rs}.

\textbf{Tablas soportadas}: \code{SigninLogs},
\code{AuditLogs}, \code{SecurityEvent}, \code{DeviceProcessEvents}
(Defender for Endpoint), \code{AzureActivity}, \code{SecurityAlert}.

\textbf{Mapeo}: depende de la tabla; el parser resuelve vía
\code{FieldConfig} propio de cada una. Ejemplo para
\code{SigninLogs}: \code{UserPrincipalName} $\to$ \code{User},
\code{IPAddress} $\to$ \code{IP}, \code{AppDisplayName} $\to$
\code{Service}, con relación \code{User -[Auth]-> Service}.

\section{CSV genérico}

Parser: \filepath{platform/parsers/src/generic_csv.rs}. Requiere un
\code{FieldConfig} explícito (producido por
\code{parser_generate} o aportado por el usuario vía la UI).

\textbf{Roles de campo}: \code{TIME}, \code{ACTOR}, \code{ACTION},
\code{TARGET}, \code{INFO}. El rol \code{TIME} es obligatorio; el
parser falla si no logra asignar uno.

\section{Generic (heurístico)}

Parser: \filepath{core/graph-engine/src/generic.rs}. Dado un
archivo de texto con líneas KV o JSON-lines, intenta:

\begin{itemize}
    \item Detectar timestamps con un catálogo de formatos comunes.
    \item Reconocer direcciones IP, dominios, hashes.
    \item Inferir relaciones por colocación de campos.
\end{itemize}

Es el fallback. Siempre produce algún grafo, pero las relaciones
pueden ser de baja calidad. La UI marca este parser con
\code{stability = experimental} en la sesión resultante.

\section{FortiAnalyzer}

Parser: \filepath{platform/parsers/src/forti.rs}. CSV propietario
con campos del estilo \code{date}, \code{time}, \code{srcip},
\code{dstip}, \code{dstport}, \code{app}, \code{action},
\code{user}.

\textbf{Mapeo}:
\begin{itemize}
    \item \code{srcip}, \code{dstip} $\to$ \code{IP}.
    \item \code{user} $\to$ \code{User}.
    \item \code{app} $\to$ \code{Service}.
\end{itemize}

\textbf{Relaciones}: \code{IP -[Connect]-> IP} con atributos
\code{dstport}, \code{app}, \code{action}.

\section{AWS Cognito}

Parser: \filepath{platform/parsers/src/cognito.rs}. JSON de
CloudWatch Logs para eventos de Cognito (signin, signup, token
refresh).

\textbf{Mapeo}: \code{username} $\to$ \code{User}, \code{sourceIp}
$\to$ \code{IP}, \code{userPoolId} $\to$ \code{Service}.

\section{IIS (W3C Extended Log)}

Parser: \filepath{platform/parsers/src/iis.rs}. Formato W3C con
\code{\#Fields} header. Maneja aliases comunes
(\code{cs-username}, \code{c-ip}, \code{cs-uri-stem}, etc.).

\textbf{Mapeo}: \code{cs-username} $\to$ \code{User}, \code{c-ip}
$\to$ \code{IP}, \code{cs-uri-stem} $\to$ \code{URL}.

\section{Tres carriles de ingesta}
\label{sec:log-three-lanes}

La ingesta productiva corre sobre el \emph{fast lane}: un
programa VRL compilado determinísticamente desde un
\code{FieldConfig} aprobado. Los parsers listados arriba son
fuentes de \code{FieldConfig}; el código que ejecuta por evento en
runtime es VRL, no el parser. Los parsers siguen activos para
ingesta legacy pre-VRL y como oráculo en
\code{mapping_regression_test}.

El detalle arquitectónico vive en el libro de Arquitectura,
capítulo \emph{Ingestor agéntico}. Esta sección cubre lo que un
cliente del sistema necesita saber para escribir contra la
superficie.

\subsection{Fast lane — VRL compilado}

Dado un \code{FieldConfig} $f$, el compilador emite un programa
VRL textual con identidad estable:

\begin{equation*}
\text{VRL}(f) = \text{compile}(f) \qquad
\text{mapping\_hash}(f) = \text{sha256}(\text{VRL}(f))
\end{equation*}

El \code{mapping_hash} se estampa en el campo
\code{provenance.mapping_hash} de cada evento canonical emitido.
Garantiza que dos eventos con el mismo hash fueron producidos por
la misma lógica de ingesta.

Contrato del compilador (\filepath{platform/vrl}):

\begin{itemize}
    \item Entrada: \code{FieldConfig} válido (semántica del
    capítulo anterior).
    \item Salida: \code{VrlProgram} (\code{struct} con el source
    VRL, el hash, y metadatos de la versión del compilador).
    \item Pura: sin I/O, sin red, sin LLM. La misma entrada produce
    idéntico output byte-a-byte.
    \item El runtime VRL ejecuta el programa sobre cada evento
    crudo y emite canonical OCSF.
\end{itemize}

\subsection{Slow lane — LLM-as-compiler}

La tool MCP \code{parser_generate} (y su endpoint HTTP
\code{POST /parser_generate}) toma una muestra del formato y
devuelve un \code{FieldConfig} propuesto. La implementación usa un
backend LLM (\filepath{platform/local-llm}) que puede ser local
(GGUF) o remoto (API). El output jamás se aplica automáticamente:
entra a la cola de review agéntica vía
\code{POST /agentic_review_list}.

Flujo típico:

\begin{enumerate}
    \item Cliente llama \code{parser_generate} con
    \code{\{sample: "...", format_hint: "fortios"\}}.
    \item Server devuelve \code{\{proposal_id, field_config,
    ocsf_mapping, confidence\}}.
    \item Humano abre \code{agentic_review_list}, aprueba o edita
    vía \code{agentic_review_approve} o \code{agentic_review_reject}.
    \item Tras aprobación, \code{publish_mapping} persiste el
    mapping en \filepath{platform/mapping-library} con
    \code{mapping_hash} estable.
    \item Próxima ingesta con ese formato usa automáticamente el
    nuevo mapping.
\end{enumerate}

\subsection{Learning lane — drift-as-signal}

Cada batch en fast lane produce métricas estructurales:

\begin{itemize}
    \item Histograma de entity types emitidos.
    \item Cardinalidad de cada campo canonical.
    \item Ratio de eventos enviados a DLQ.
    \item Distribución de timestamps observados.
\end{itemize}

El endpoint \code{POST /schema_drift_detect} (tool MCP
\code{schema_drift_detect}) toma un batch y devuelve un score de
drift por campo usando Kolmogorov-Smirnov contra el perfil del
mapping actual. Si el drift supera tolerancia, el sistema:

\begin{enumerate}
    \item \textbf{No} para la ingesta — sigue usando el mapping
    actual.
    \item Abre propuesta de re-review en la cola agéntica.
    \item El slow lane procesa la propuesta; humano decide si
    aceptar la evolución.
\end{enumerate}

El contrato es explícito: drift no es error, es señal. El sistema
nunca deja de ingestar por descubrir un cambio aguas arriba.

\section{Canonical OCSF}
\label{sec:log-ocsf}

Todo mapping emite, como salida canonical, eventos con shape
\textbf{OCSF v1.4} (Open Cybersecurity Schema Framework). Campos
principales garantizados:

\begin{table}[H]
\centering
\small
\begin{tabular}{@{}ll@{}}
\toprule
Campo & Semántica \\
\midrule
\code{time}                     & timestamp ISO-8601 UTC \\
\code{category_uid}             & categoría OCSF (ej. 4 Network) \\
\code{class_uid}                & clase específica \\
\code{activity_id}              & acción normalizada \\
\code{src_endpoint.ip}          & origen IP \\
\code{dst_endpoint.ip}          & destino IP \\
\code{actor.user.name}          & principal \\
\code{device.hostname}          & host \\
\code{observables[]}            & entidades extraídas para el grafo \\
\code{provenance.mapping_hash}  & hash del VRL que produjo este evento \\
\code{provenance.source_format} & formato raw de origen \\
\bottomrule
\end{tabular}
\caption{Subconjunto de campos OCSF garantizados.}
\end{table}

Cualquier campo adicional específico del formato raw se emite bajo
\code{raw.\{...\}} como escape hatch: el grafo no lo usa pero el
analista puede consultarlo.

\section{FieldConfig como contrato}

Todos los parsers convergen a un \code{FieldConfig} que describe
roles y mapeos. Este FieldConfig se puede:

\begin{itemize}
    \item \textbf{Inferir} vía \code{parser_generate} (el sistema
    propone, el usuario confirma).
    \item \textbf{Compartir} vía \code{publish_mapping} al
    \code{mapping_library}.
    \item \textbf{Versionar}: cada mapping tiene su versión
    propia; cambios requieren nueva versión.
    \item \textbf{Compilar a VRL} vía \code{platform/vrl} para
    portarlo a pipelines externos (Vector, Fluent Bit).
\end{itemize}

\section{Cómo agregar un parser}

Resumen (el detalle vive en el libro de Arquitectura, capítulo
\emph{Capa Platform}):

\begin{enumerate}
    \item Crear el archivo en \filepath{platform/parsers/src/<nombre>.rs}.
    \item Implementar \code{LogParser}.
    \item Registrar en \filepath{platform/parsers/src/lib.rs}.
    \item Agregar fixtures a \filepath{platform/parsers/tests/}.
    \item Agregar regla de detección a \code{format_registry}.
\end{enumerate}

No es necesario tocar ninguna otra capa.
